%% =============================================================================
%% chapters/07-presentation-iris.tex
%% PART II: Platform Capabilities and Modules
%% Chapter 7: Presentation Services & User Interfaces (Iris BEFE, Clinical, Console, Administration)
%% =============================================================================

\chapter{Presentation Services \& User Interfaces (\texttt{Iris})}
\label{chap:presentation_iris}

The Iris subsystem is Harmonia's presentation tier, providing decoupled, modern web user interfaces and a high-performance Backend-For-Frontend (BEFE) gateway.

\begin{figure}[htbp]
    \centering
    \begin{adjustbox}{max width=\textwidth}
        %% =============================================================================
%% diagrams/fig-iris-architecture.tex
%% ArchiMate 3.2 Iris Presentation Tier & BEFE Ecosystem Architecture View
%% =============================================================================
\begin{tikzpicture}[node distance=1.0cm and 0.8cm, auto]

    % --- Tier 1: Iris Presentation SPAs (Port 3000, 3001, 3002) ---
    \node[archi-app-component] (ui_clinical) at (1.5, 9.5) {%
        \archibadge{App Component}\archiname{Iris Clinical UI}\\\archidesc{Vue 3 SPA / Port 3000}%
    };

    \node[archi-app-component] (ui_admin) at (7.0, 9.5) {%
        \archibadge{App Component}\archiname{Iris Administration}\\\archidesc{Vue 3 SPA / Port 3002}%
    };

    \node[archi-app-component] (ui_console) at (12.5, 9.5) {%
        \archibadge{App Component}\archiname{Iris Console UI}\\\archidesc{Vue 3 SPA / Port 3001}%
    };

    % --- Tier 2: Iris BEFE Gateway & Security (Ports 8080, 8090) ---
    \node[archi-app-interface] (befe_fhir) at (4.25, 7.0) {%
        \archibadge{App Interface}\archiname{Clinical FHIR Gateway}\\\archidesc{JAX-RS REST / Port 8080}%
    };

    \node[archi-app-interface] (befe_ops) at (12.5, 7.0) {%
        \archibadge{App Interface}\archiname{Operations Server}\\\archidesc{OperationsServer / Port 8090}%
    };

    \node[archi-app-component] (themis_auth) at (4.25, 4.8) {%
        \archibadge{Security Component}\archiname{Themis Security}\\\archidesc{RBAC Policy Evaluator}%
    };

    % --- Tier 3: Messaging, Workflow & Persistence ---
    \node[archi-app-component] (petasos_queue) at (1.5, 2.5) {%
        \archibadge{App Component}\archiname{Petasos Broker}\\\archidesc{ActiveMQ Artemis / 61616}%
    };

    \node[archi-app-component] (ponos_engine) at (12.5, 4.8) {%
        \archibadge{App Component}\archiname{Ponos WorkEngine}\\\archidesc{Praxis Engine / Port 8083}%
    };

    \node[archi-node] (mneme_grid) at (7.0, 4.8) {%
        \archibadge{Technology Node}\archiname{Mneme Cache Grid}\\\archidesc{Infinispan / Port 11222}%
    };

    \node[archi-data-object] (mnemosyne_db) at (7.0, 2.5) {%
        \archibadge{Data Object}\archiname{Mnemosyne DB}\\\archidesc{PostgreSQL FHIR Store}%
    };

    % --- Relationships & Interactions ---
    % UI to BEFE Gateway
    \draw[archi-flow] (ui_clinical) -- node[left, font=\tiny] {FHIR R5 / REST} (befe_fhir);
    \draw[archi-flow] (ui_admin) -- node[right, font=\tiny] {FHIR CRUD / Changes} (befe_fhir);
    \draw[archi-flow] (ui_console) -- node[right, font=\tiny] {Status / Reload} (befe_ops);

    % Security Enforcement
    \draw[archi-serving] (themis_auth) -- node[right, font=\tiny] {Evaluate Auth} (befe_fhir);

    % Clinical Gateway to Infrastructure
    \draw[archi-flow] (befe_fhir) -- node[above, font=\tiny] {HotRod Client} (mneme_grid);
    \draw[archi-flow] (befe_fhir) -- node[left, font=\tiny] {JMS Change Pragma} (petasos_queue);

    % Operations Server to Workflow & Cache
    \draw[archi-serving] (befe_ops) -- node[right, font=\tiny] {Engine Telemetry} (ponos_engine);
    \draw[archi-access] (befe_ops) -- node[above, font=\tiny] {Cache Metrics} (mneme_grid);

    % Workflow & Persistence
    \draw[archi-flow] (petasos_queue) -- node[below, font=\tiny] {Async Dispatch} (ponos_engine);
    \draw[archi-flow] (mneme_grid) -- node[right, font=\tiny] {Write-Behind} (mnemosyne_db);

\end{tikzpicture}

    \end{adjustbox}
    \caption{Iris Presentation Tier and BEFE Gateway Architecture}
    \label{fig:iris_presentation_tier}
\end{figure}

\section{Presentation Single-Page Applications (SPAs)}

Harmonia provides three purpose-built Vue 3 / TypeScript web applications served via Nginx:

\subsection{1. \texttt{iris-clinical} (Clinical Portal)}
The clinical portal provides care teams with fast, intuitive access to longitudinal patient histories, encounter summaries, lab observations, and diagnostic reports.

\subsection{2. \texttt{iris-console} (Operations \& Telemetry Portal)}
The operations console provides Site Reliability Engineers (SREs), system administrators, and integration engineers with comprehensive operational telemetry, message flow tracing, and incident remediation capabilities.

\begin{archinote}
\textbf{Terminology Reconciliation}: Legacy documentation previously referred to this portal as \texttt{iris-monitor}. The canonical, repository-backed identifier is standardized across all modules and documentation as \texttt{iris-console}.
\end{archinote}

The console is structured around five primary operational perspectives:
\begin{enumerate}
    \item \textbf{SUBSYSTEMS (\texttt{/subsystems})}: Answers \textit{"Is Harmonia healthy?"} Visualizes platform-wide health states across all nine core subsystems (Pylai, Petasos, Energeia [Ponos, Praxis], Mneme, Mnemosyne, Calliope, Themis, Agora, Iris). Provides domain-neutral runtime instance tables with slide-out Pod/container detail drawers, downstream dependency health ratios with round-trip latencies, and native SVG micro-chart trend statistics across configurable time windows (15m, 1h, 6h, 24h).
    \item \textbf{QUEUES (\texttt{/queues})}: Answers \textit{"Is work moving?"} Monitors Apache ActiveMQ Artemis message broker topology, queue depths, enqueue and dequeue rates, consumer/producer ratios, oldest message age, message redeliveries, and Dead-Letter Queue (DLQ) accumulation surveillance with historical depth sparklines.
    \item \textbf{WORKFLOWS (\texttt{/workflows})}: Answers \textit{"What work is Harmonia currently performing?"} Displays registered Praxis execution sequences, active, queued, completed, and failed task volumes, processing throughput rates, P95 durations, and deep diagnostic drill-down into individual Pragma execution envelopes and Ergon checkpoint timelines.
    \item \textbf{EVENTS (\texttt{/events})}: Answers \textit{"What happened to a particular interaction?"} Cross-subsystem end-to-end diagnostic sequence tracing across protocol gateways, message queues, workflow processors, and persistence. Supports multi-attribute filtering (Correlation ID, Causation ID, Message ID, Pragma ID, Subsystem, Event Type, Status, Time Range) rendered as an operational vertical flowchart timeline.
    \item \textbf{ALERTS (\texttt{/alerts})}: Answers \textit{"What currently requires operator attention?"} Actionable operational conditions categorized by severity (Critical, Warning, Information) and lifecycle status (Active, Acknowledged, Resolved) accompanied by concrete operator remediation guidance and audit-governed acknowledgment actions.
\end{enumerate}

\begin{principlebox}{Zero-PHI Presentation Boundary (Invariant 7)}
In strict compliance with Invariant 7, the Iris Operations Console and Operations Backend enforce a complete Zero-PHI boundary. All diagnostic event timelines, queue inspections, and workflow details display only technical identifiers (\texttt{correlationId}, \texttt{pragmaId}, \texttt{messageId}, durations, and error reason codes). Clinical payloads, FHIR resource bodies, HL7 segments, and patient identifiers are strictly excluded from UI rendering and operations API responses.
\end{principlebox}

\subsection{3. \texttt{iris-administration} (Directory Administration Portal)}
Provides directory stewards with UI forms for managing practitioners, roles, organizations, and endpoints.

\begin{principlebox}{Iris-Administration Role Invariant}
In accordance with Invariant 3, \texttt{iris-administration} is strictly a frontend presentation client consuming FHIR REST APIs. It contains zero database drivers, no JPA/Hibernate entities, and does not execute Provider Registry state machines or referential integrity verification. All authoritative registry governance is executed server-side.
\end{principlebox}

\section{Backend-For-Frontend Gateway (\texttt{iris-befe})}

\texttt{iris-befe} is a WildFly 31.0.1 Jakarta EE microservice operating on dual HTTP ports:
\begin{itemize}
    \item \textbf{Port 8080 (Clinical Gateway)}: Serves authenticated clinical FHIR REST endpoints to \texttt{iris-clinical}.
    \item \textbf{Port 8090 (Operations Gateway)}: Serves system telemetry, normalized subsystem models, queue statistics, workflow diagnostics, and alert remediation endpoints to \texttt{iris-console}.
\end{itemize}

\subsection{Operations Backend Architecture (Port 8090)}
The operations presentation tier communicates exclusively through \texttt{iris-befe} port 8090 via the \texttt{OperationsServerManager}. Telemetry collection is orchestrated by the \texttt{OperationsAggregatorService} through a decoupled provider architecture:
\begin{itemize}
    \item \textbf{Subsystem Health Provider SPI (\texttt{SubsystemHealthProvider})}: Concrete provider implementations for each core subsystem (\texttt{PylaiHealthProvider}, \texttt{PetasosHealthProvider}, \texttt{EnergeiaHealthProvider}, \texttt{MnemeHealthProvider}, \texttt{MnemosyneHealthProvider}, \texttt{CalliopeHealthProvider}, \texttt{ThemisHealthProvider}, \texttt{AgoraHealthProvider}, and \texttt{IrisHealthProvider}) normalize disparate telemetry sources into uniform operational models.
    \item \textbf{Hybrid Telemetry Sourcing}: State is harvested primarily from Infinispan clustered caches (\texttt{modulestatus-cache}, \texttt{task-cache}, \texttt{messagequeue-cache}) supplemented by non-blocking subsystem actuator health probes and ActiveMQ Artemis broker management APIs.
    \item \textbf{Kubernetes Runtime Discovery (\texttt{KubernetesInstanceProvider})}: Dynamically discovers live Pods in the \texttt{harmonia} namespace using in-cluster ServiceAccount credentials, mapping Pod lifecycle and resource metrics to domain-neutral instances with graceful fallback to Docker Compose.
    \item \textbf{Themis Default-Deny Governance}: Operations endpoints enforce strict RBAC checks using \texttt{themis-api} contracts, requiring \texttt{operations.viewer} or \texttt{operations.admin} authorities.
    \item \textbf{Partial Failure Resilience}: When downstream telemetry providers experience timeouts or temporary unavailability, the aggregator tags the corresponding subsystem as \texttt{UNKNOWN} or \texttt{UNAVAILABLE} with stale indicators, guaranteeing that platform summaries continue responding within bounded SLAs ($< 200$\,ms).
\end{itemize}

\begin{figure}[htbp]
    \centering
    \begin{adjustbox}{max width=\textwidth}
        %% =============================================================================
%% diagrams/fig-iris-federation.tex
%% Iris Presentation Federation (SPAs, BEFE Gateway, Hot Rod & REST Endpoints)
%% =============================================================================
\begin{tikzpicture}[node distance=1.0cm and 0.8cm, auto]

    % --- Background Plates ---
    \draw[archi-plate-bg] (-0.8, 6.0) rectangle (15.8, 9.4);
    \node[archi-plate-title] at (-0.6, 9.2) {Iris Presentation SPAs (Vue 3 / TypeScript / Pinia)};

    \draw[archi-plate-bg] (-0.8, 0.4) rectangle (15.8, 5.2);
    \node[archi-plate-title] at (-0.6, 5.0) {Backend-For-Frontend (BEFE) Gateway \& Downstream Services};

    % --- SPAs ---
    \node[archi-app-component, minimum width=3.4cm] (ui_clin)  at (1.5, 7.8) {\archibadge{SPA}\archiname{iris-clinical}\\\archidesc{Clinical Port 80/3000}};
    \node[archi-app-component, minimum width=3.4cm] (ui_cons)  at (6.5, 7.8) {\archibadge{SPA}\archiname{iris-console}\\\archidesc{Console Port 80/3001}};
    \node[archi-app-component, minimum width=3.4cm] (ui_admin) at (11.5, 7.8) {\archibadge{SPA}\archiname{iris-administration}\\\archidesc{Admin Port 80/3002}};

    % --- BEFE & Services ---
    \node[archi-app-component, minimum width=4.5cm] (befe)      at (4.0, 3.4) {\archibadge{BEFE Gateway}\archiname{iris-befe}\\\archidesc{WildFly 31 / JAX-RS / CDI}};
    \node[archi-app-component, minimum width=4.5cm] (pylai_reg) at (11.5, 3.4) {\archibadge{REST Gateway}\archiname{pylai-fhir-registry}\\\archidesc{Spring Boot FHIR R5}};

    \node[archi-system-software, minimum width=4.5cm] (mneme) at (4.0, 1.4) {\archibadge{Cache}\archiname{mneme-cluster}\\\archidesc{Infinispan Hot Rod 11222}};
    \node[archi-data-object, minimum width=4.5cm] (mnemo)     at (11.5, 1.4) {\archibadge{Storage}\archiname{mnemosyne-clinical}\\\archidesc{PostgreSQL JPA R5}};

    % --- Flows ---
    \draw[archi-flow] (ui_clin) -- node[archi-lbl]{REST} (befe);
    \draw[archi-flow] (ui_cons) -- node[archi-lbl]{REST / WS} (befe);
    \draw[archi-flow] (ui_admin) -- node[archi-lbl]{FHIR REST} (pylai_reg);

    \draw[archi-serving] (mneme) -- node[archi-lbl]{Hot Rod} (befe);
    \draw[archi-flow] (pylai_reg) -- node[archi-lbl]{JPA} (mnemo);

\end{tikzpicture}

    \end{adjustbox}
    \caption{Iris UI Federation, Hot Rod Caching \& REST Endpoint Integration}
    \label{fig:iris_federation}
\end{figure}

\newpage
